\section{Introduction}
\label{sec:introduction}

The pluricanonical systems of an algebraic fiber space reflect both
the geometry of its fibers and the positivity available from its base.
When the geometric generic fiber has Kodaira dimension zero, the
easy-addition bound says that the total space has Kodaira dimension
at most the dimension of the base. The zero-Kodaira case of Schnell's
question asks whether a numerical comparison with an ample divisor on
the base forces equality.
Here an \emph{algebraic fiber space} means a surjective morphism with
connected fibers between projective varieties.

We approach this question through a good canonical model of the total
space. The geometric argument below shows what such a model supplies;
the smooth canonical good-model theorem of
\cite[Corollary~11.2]{LogAbundance2026} then provides the model.
Its nonvanishing consequence also permits Schnell's reduction to pass
from zero-Kodaira fibers to the Campana--Peternell inequality and the
general fiber-space conclusion stated below.

The formulation in \cite[Conjecture~1.2]{Kim2025} compares
a positive multiple of $K_X$ with the pullback of an ample Cartier divisor on the base.
The difficulty is that pseudo-effectivity concerns a numerical divisor
class, whereas Kodaira dimension measures actual sections. The
following theorem gives the positive conclusion in this formulation;
its assertion has also been proved by Zou, as discussed below.

\begin{theorem}[Schnell's zero-Kodaira fiber-space conclusion]
\label{thm:schnell}
Let $f\colon X\to Y$ be a surjective morphism with connected fibers
between smooth connected projective complex varieties. Let $F$ be its
smooth geometric generic fiber and assume $\kappa(F)=0$. Suppose that
there are an ample Cartier divisor $H$ on $Y$ and a positive integer
$m_0$ such that $m_0K_X-f^*H$ is pseudo-effective. Then
\[
 \kappa(X)=\dim Y.
\]
\end{theorem}

The ample comparison is essential. For example, let $F$ be smooth,
connected and projective with torsion canonical bundle. The projection
$F\times\mathbb P^1\to\mathbb P^1$ has fibers of Kodaira dimension
zero, but its total space has Kodaira dimension $-\infty$.
The class $m_0K_{F\times\mathbb P^1}-f^*H$ has negative degree on
the moving curves $\{x\}\times\mathbb P^1$ for every ample $H$;
it therefore cannot be pseudo-effective: for any fixed effective
divisor, a general member of this family is not contained in its
support and has nonnegative intersection with it. In contrast, for a
smooth connected projective $Y$ with $K_Y$ ample, the projection
$F\times Y\to Y$ satisfies the numerical
hypothesis with $H=K_Y$ for every $m_0\geq1$, since the compared
class is numerically $(m_0-1)f^*K_Y$. Its Kodaira dimension is
$\dim Y$.

The canonical good-model theorem used below also supplies canonical
nonvanishing in every dimension. With this input, Schnell's reduction
\cite[Sections~4--10]{Schnell2022} promotes
Theorem~\ref{thm:schnell} to the Kodaira-dimension form of the
Campana--Peternell conjecture.

\begin{corollary}[Campana--Peternell]
\label{cor:campana-peternell}
Let $X$ be a smooth connected projective complex variety, let $D$ be
an effective Cartier divisor on $X$, and let $m_0$ be a positive
integer. If $m_0K_X-D$ is pseudo-effective, then
\[
 \kappa(X)\geq\kappa(D).
\]
\end{corollary}

The same reduction gives the general fiber-space conclusion, including
the existence of sections after subtracting the original ample
pullback.

\begin{corollary}[Schnell's general fiber-space conclusion]
\label{cor:schnell-general}
Let $f\colon X\to Y$ be a surjective morphism with connected fibers
between smooth connected projective complex varieties, and let $F$ be
a very general smooth fiber. Suppose that $H$ is an ample Cartier
divisor on $Y$ and $m_0$ is a positive integer such that
$m_0K_X-f^*H$ is pseudo-effective. Then
\[
 \kappa(X)=\kappa(F)+\dim Y.
\]
Moreover, there are positive integers $r$ and $\ell_0$ such that
\[
 H^0\bigl(X,\cO_X(\ell rK_X-f^*H)\bigr)\ne0
 \qquad\text{for every integer }\ell\geq\ell_0.
\]
\end{corollary}

The last assertion is effectivity for all sufficiently large and
divisible pluricanonical degrees, with no uniform choice of $r$ or
$\ell_0$ asserted. The numerical hypothesis itself implies
$\kappa(F)\geq0$ by restriction and canonical nonvanishing; it is not
replaced by that weaker fiber condition. Corollary~\ref{cor:campana-peternell}
also applies to an effective rational divisor after clearing
denominators, as explained in Section~\ref{sec:general-consequences}.

\subsection{The problem and previous approaches}

Campana and Peternell formulated the comparison
$NK_X=A+B$, with $N$ a positive integer, $A$ effective and $B$
pseudo-effective, in their
study of positivity of the cotangent bundle
\cite[Conjecture~2.4]{CP2011}. The proposed inequality
$\kappa(X)\geq\kappa(A)$ includes canonical nonvanishing when
$A=0$: pseudo-effective $K_X$ should have a nonzero pluricanonical
section. For general $A$, the conclusion also requires the canonical
systems to have at least the image dimension supplied by $A$.
When $\kappa(X)\geq0$, Campana and Peternell reduce this comparison
to the general fibers of the Iitaka fibration of $X$, which have
Kodaira dimension zero \cite[Proposition~2.6]{CP2011}. In that
zero-Kodaira setting, a good minimal model has torsion canonical
class; their pushforward argument then forces $\kappa(A)=0$
\cite[Proposition~2.7]{CP2011}. This already exhibits how a good
model can convert numerical information into a statement about
sections.

Schnell developed this circle of questions in his study of singular
metrics, nonvanishing and the Campana--Peternell conjecture
\cite[Sections~4--10]{Schnell2022}. His Conjecture~10.1 allows fibers
of nonnegative Kodaira dimension and asks for nonvanishing after
subtracting an ample pullback. Under canonical nonvanishing, his
Section~9 reduces the conjecture to fibers of Kodaira dimension zero;
Lemma~7.1 supplies an effective pluricanonical divisor after a positive
base twist. If the fiber has positive Kodaira dimension, the Iitaka
fibration of that divisor has a strictly larger base, while retaining
the numerical comparison. Iteration reaches the zero-Kodaira case.
Section~8 connects the resulting equality to effectivity after
subtracting the original ample pullback. This uses the Iitaka-addition
criterion of Fujita and Mori, in the form recorded in
\cite[Lemma~4.6]{PopaSchnell2014}. He proves a special case when the
canonical divisor of the base is pseudo-effective
\cite[Theorem~12.1]{Schnell2022}.

Kim applies the canonical bundle formula to the same fiber-space
problem. His Theorem~1.3 proves the conclusion when
$K_Y+(1-\varepsilon)B_Y$ is pseudo-effective for some
$\varepsilon>0$, where $B_Y$ is the discriminant on the birational
setup of that theorem, or when the canonical class of the general
fiber is rigid \cite{Kim2025}. Here rigidity means uniqueness of
the closed positive $(1,1)$-current representing that class.
Kim also recalls a known algebraic proof under the existence of a
good minimal model of the general fiber
\cite[Section~1.3]{Kim2025}. Zou obtains the same
conclusion as Theorem~\ref{thm:schnell}, without these additional
hypotheses, by a canonical-bundle-formula argument
\cite[Theorems~1.3 and~5.1]{Zou2025}.

We give a direct mixed-intersection proof of
Theorem~\ref{thm:schnell} from a good minimal model of the
\emph{total space}. We separate that geometric argument from the
model-existence theorem that it uses. The latter is the smooth
canonical good-model theorem of
\cite[Corollary~11.2]{LogAbundance2026}, quoted in
Theorem~\ref{thm:canonical-good-models}. Its nonvanishing consequence
and Schnell's published reduction then yield
Corollaries~\ref{cor:campana-peternell} and~\ref{cor:schnell-general}.

\subsection{How the numerical condition produces the lower bound}

A good klt minimal model of $X$ consists of a normal projective
$\Q$-factorial klt variety $V$ with semiample canonical divisor and a
$K_X$-negative birational contraction $X\dashrightarrow V$.
For compatible canonical divisors, the comparison on a smooth common
resolution $X\xleftarrow{p}W\xrightarrow{q}V$ has the form
\begin{equation}\label{eq:intro-comparison}
 p^*K_X=q^*K_V+E,\qquad E\geq0,\qquad q_*E=0.
\end{equation}
The last condition means that every component of $E$ is exceptional
over $V$. The model therefore provides sections and also specifies
the error in transferring them to $X$.

Choose a globally generated Cartier multiple $rK_V$. The comparison
makes $rE$ an effective Cartier divisor. Pulling back sections of
$rK_V$ and multiplying by the canonical section of $rE$ gives a
pluricanonical subsystem on $X$ with the same image dimension.
It remains to force that dimension to be at least $\dim Y$.

If the image were smaller, Lemma~\ref{lem:mixed-positivity} would
provide a mixed complete-intersection curve class $C$ on $W$ with
three properties: it pairs nonnegatively with every pseudo-effective
divisor, it annihilates both $q^*K_V$ and $E$, and it pairs positively
with $(fp)^*H$. Its divisor factors are pulled back from $V$, so
the exceptional vanishing follows from the Cartier projection formula.
The strict positivity uses a cotangent direction from the map to $Y$
beyond those supplied by the canonical morphism. Thus
\[
 (m_0p^*K_X-(fp)^*H)\cdot C=-(fp)^*H\cdot C<0,
\]
contrary to the pulled-back numerical hypothesis. Section~\ref{sec:geometry}
constructs the class and proves each of these properties.

This argument uses the classical intersection theory of Cartier
divisors and proper pushforward \cite[Chapter~2]{Fulton1998}.
In particular, the maps to $Y$ and to the canonical image need no factorization
relation. Section~\ref{sec:geometry} gives the complete geometric
proof, including the upper bound and the transfer of sections.
Section~\ref{sec:good-model-interface} states the exact good-model
input, verifies its hypotheses and completes
Theorem~\ref{thm:schnell}. Section~\ref{sec:general-consequences}
explains the nonvanishing and fiber interfaces, then proves the
Campana--Peternell and general Schnell conclusions.

\subsection{Conventions}

All varieties are integral and projective over $\C$ unless specified
otherwise. Smoothness of an algebraic fiber space refers here to its
source and target; special fibers of the morphism may be singular.
Canonical divisors on birational models are chosen compatibly.
For a rational Cartier divisor $L$, the Iitaka dimension $\kappa(L)$
is the maximum dimension of the rational images of its nonempty
complete systems $|mL|$, with $m>0$ clearing the Cartier index.
It is $-\infty$ if all these systems are empty. We write
$\kappa(X)=\kappa(K_X)$.

We write $\PE(X)$ for the closure of the cone of effective real
divisor classes in $N^1(X)_\R$. A divisor is pseudo-effective when
its numerical class belongs to this cone. A rational Cartier divisor
is semiample when a positive Cartier multiple is globally generated.
The section constructions below use this actual line-bundle property.
